\section{The Frame Attribution Problem}
\label{sec:problem}

This section describes the gap between what hitch-capture tools observe and
what an engineer needs (\autoref{sec:gap}), the limitations of existing
approaches (\autoref{sec:limits}), what an in-kernel declarative tracer can
express (\autoref{sec:ftrace}), and the scope we assume (\autoref{sec:scope}).

\subsection{Challenge: The Symptom--Cause Gap}
\label{sec:gap}

The cheapest signal of a hitch is its symptom, the frame time. It is known
only when the frame ends, and it says nothing about cause. The cause lies in
events that happened during the frame, on the presenting thread and on the
threads it waited for: when each thread was off-CPU, why, and who ended each
wait. We call the distance between the two the \emph{symptom--cause gap}.
How a tool bridges it determines two trade-offs.

\paragraph{Trade-off 1: Output volume versus detail}
A tool can keep the raw events and decide later, as a flight recorder does,
paying a full window per hitch; or it can reduce events to a summary as they
happen, paying little but answering only the questions built into the
summary.

\paragraph{Trade-off 2: Gate timing versus sensitivity}
Whether a stall matters depends on whether its frame ends up over budget,
which is unknown while the stall is happening. A gate evaluated per stall
must use a per-stall threshold and misses frames made of many small stalls;
a gate evaluated at frame end must carry state across the frame. Consider a
16.7\,ms frame with 7\,ms of work and forty waits of 200--600\,$\mu$s:
futex hand-offs, minor faults, two short preemptions. The frame takes
23\,ms. No single stall is remarkable, and normal frames contain dozens of
similar stalls; what distinguishes the bad frame is the sum. By the time the
thirty-fifth stall shows the frame going bad, a per-event filter has already
dropped the first thirty-four.

\subsection{Limitations of Existing Approaches}
\label{sec:limits}

\autoref{tab:compare} summarizes the approaches.

\begin{table*}[t]
\centering
\footnotesize
\setlength{\tabcolsep}{4pt}
\begin{tabular}{@{}lcccccc@{}}
\toprule
 & \textbf{Gate} & \textbf{Sub-threshold} & \textbf{Output} & \textbf{Cause} & \textbf{Cross-thread} & \textbf{Custom} \\
 & \textbf{evaluated} & \textbf{sums} & \textbf{per hitch} & \textbf{attribution} & \textbf{chain} & \textbf{kernel code} \\
\midrule
Flight recorder~\cite{ftracehist,perfetto-triggers,desfossez2016latency} & frame end, symptom & raw events & whole window & manual & manual & none / module \\
Per-event threshold~\cite{perf-offcpu-thresh,gregg2019bpf} & each stall & \xmark & per stall & stack & one hop (offwake) & BPF \\
Delay accounting~\cite{delayacct,psi} & polled & \cmark & counters & categories, no futex & \xmark & none \\
Per-request accounting~\cite{varmri} & request end & \cmark & one record & categories & \xmark & kernel patch \\
Offline critical path~\cite{giraldeau2016wait,zhou2018wperf,perfetto-critpath} & offline & \cmark & full trace & \cmark & \cmark & none \\
\textbf{\sys} & \textbf{frame end, cause} & \cmark & \textbf{one record} & \cmark & \textbf{online, bounded} & \textbf{verified BPF} \\
\bottomrule
\end{tabular}

\caption{Comparison of approaches to capturing and explaining frame hitches.}
\label{tab:compare}
\end{table*}

\subsubsection{Flight Recorders}

A flight recorder records continuously into an overwriting buffer and
freezes it on a trigger. The frame-time trigger itself needs no custom
kernel code: an ftrace uprobe event on the present call, a histogram trigger
that computes the delta between consecutive presents into a synthetic event,
and a \texttt{snapshot} trigger filtered on that delta implement ``freeze the
buffer when a frame exceeds 16.6\,ms'' entirely in the kernel~\cite{ftracehist}.
Perfetto does the same with a user-space trigger~\cite{perfetto-triggers}, and
gpuvis~\cite{gpuvis} with a manual one. The costs are those listed in the
introduction: the trigger is symptom-only, every trigger dumps the window,
and diagnosis is manual. In steady state, however, a flight recorder moves no
data at all, so any claim that an eBPF tool ``reduces data movement'' must be
measured against it rather than against streaming tracers (\autoref{sec:eval-overhead}).

\subsubsection{Per-Event Thresholds}

\texttt{perf record --off-cpu} aggregates off-CPU time in a BPF map and,
with \texttt{--off-cpu-thresh}, emits a sample for each off-CPU interval
longer than a threshold~\cite{perf-offcpu-thresh}; bcc's
\texttt{offwaketime} pairs each blocked stack with its waker's stack, and
\texttt{runqslower} prints run-queue delays above a
threshold~\cite{gregg2019bpf}. These tools summarize in the kernel, but the
gate applies to each stall, not to the frame, so the forty-stall frame of
\autoref{sec:gap} produces nothing at a 1\,ms threshold and a stream
indistinguishable from normal frames at 100\,$\mu$s.

\subsubsection{Delay and Per-Request Accounting}

The kernel's delay accounting already accumulates, per task and without a
threshold, time spent waiting for a CPU, for block I/O, in swap-in, in
reclaim, in thrashing and in compaction~\cite{delayacct}. Reading the
counters at consecutive presents yields per-frame sums with no kernel code,
which makes it the baseline any frame-scoped tool must beat. It lacks a
futex category (futex waits are plain sleeps), any notion of the waker, and
per-frame scoping without diffing every thread from user space every frame.
VarMRI applies the same accumulate-and-emit pattern per server request, and
had to patch the kernel after finding the required hooks unavailable to
eBPF~\cite{varmri}.

\subsubsection{Offline Critical Path}

Critical-path analysis~\cite{giraldeau2016wait,zhou2018wperf,perfetto-critpath}
answers the chain question precisely, but over a full trace after the fact.
Its algorithm is not our contribution. Our contribution is executing a
bounded form of it online, in the kernel, gated on frame outcome.

\subsection{What an In-Kernel Declarative Tracer Can Express}
\label{sec:ftrace}

Since ftrace can decide in the kernel which events to keep, the question is
not \emph{whether} the kernel should decide but \emph{what} the decision can
depend on. ftrace's in-kernel language is intentionally declarative: stateless
filters over an event's fields, triggers, and histograms whose variables
pair a start event with an end event by key~\cite{ftracehist}. Running
arbitrary logic at a tracepoint is the problem eBPF's verifier exists to
make safe. \autoref{tab:ftrace} lists the pieces of frame attribution and
which side of that line each falls on.

\begin{table}[t]
\centering
\small
\begin{tabular}{@{}p{5.2cm}cc@{}}
\toprule
\textbf{Logic} & \textbf{ftrace} & \textbf{eBPF} \\
\midrule
Frame-time predicate at present & \cmark & \cmark \\
Per-stall duration threshold & \cmark & \cmark \\
Per-thread sum reset at each frame boundary & \xmark$^\dagger$ & \cmark \\
Sum over all of a process's threads at present & \xmark & \cmark \\
Classify by task state not in the event (\texttt{in\_iowait}, syscall in progress, preemption count) & \xmark & \cmark \\
Multi-hop wake chain with interval overlap & \xmark & \cmark \\
Commit or discard a computed record at frame end & \xmark$^\ddagger$ & \cmark \\
\bottomrule
\end{tabular}

\caption{Parts of frame attribution expressible in ftrace's in-kernel
filter/trigger/histogram language. $^\dagger$Histogram variables pair two
events by key; a running sum reset by a third event is outside the documented
model \todo{verify on 7.0}. $^\ddagger$Only whole-buffer snapshots.}
\label{tab:ftrace}
\end{table}

The first two rows show that ftrace can already make the \emph{detection}
record proportional to hitches: a filtered synthetic event is tens of bytes.
The remaining rows show that it cannot make the \emph{attribution} record
proportional to hitches, because attribution needs state accumulated across
the frame and a decision taken at its end. That is the part \sys moves into
eBPF, and it is the claim we evaluate. We do not claim that eBPF reduces
per-event cost: a BPF program on \texttt{sched\_switch} costs at least as
much as an ftrace ring-buffer write, and whether frame-scoped reduction
pays for itself is an empirical question (\autoref{sec:eval-overhead}).

\subsection{Scope and Assumptions}
\label{sec:scope}

\paragraph{Target} We assume an unmodified game, native or under Proton, that
presents through Vulkan, on a stock kernel with BTF, traced by an engineer
with \texttt{CAP\_BPF} and \texttt{CAP\_PERFMON}. \sys is a bench
instrument for platform, scheduler and driver engineers, not a consumer tool.

\paragraph{In scope} Time the presenting thread spends off-CPU or inside kernel
stall paths: preemption, run-queue delay, futex, poll, block I/O, sleeps,
GPU fence waits, present and display pacing, page faults, reclaim and
compaction, and, through the chain, the same on any thread or process it
waited for.

\paragraph{Out of scope} What the game's own code does on the CPU, and what the
GPU does internally. \sys reports on-CPU time so that a user can tell when a
hitch is \emph{not} kernel-observable and switch tools; that negative answer
is itself useful (\autoref{sec:eval-gate}). Shader compilation, the most
widely reported source of PC stutter, is on-CPU work and outside our scope.
